% ---
% id: protocol-compiler/prop:conservative
% kind: prop
% title: "Dropping a trusted marginal enlarges the feasible set"
% summary: >-
%   Dropping a trusted marginal and reading homogeneous constraints as scalars
%   enlarges the feasible set, for every admissible marginal, so it can only
%   lower a minimised objective. It compares the two source modes and is no
%   rule of emission.
% prominence: standard
% status: accepted
% uses: [lem:hermitian-constraint-values, def:homog, def:feasible, def:exact]
% mentions: [rem:notfallback, transmitter/rem:da1]
% assumes: []
% refs: []
% provenance: newly-derived
% lean:
% ---

A trusted marginal is one more condition on the states the adversary may
hold. Dropping it, and reading each homogeneous entry as its scalar, can only
admit more states; the proposition states this as an inclusion of
\tref[feasible set]{feasible sets}, and Remark~\ref{rem:notfallback} says what it licenses and what it does
not.

\begin{proposition}\label{prop:conservative}
Let a specification request trusted-source mode with marginal $\tau_A$
satisfying D(ii), and a licensed list
$\{(\hat O_k,\bowtie_k,c_k)\}_{k=1}^{m}$, $m\ge0$, satisfying D(iii) --- every
$\hat O_k$ a \tref{block list} against $\Smap$ and Hermitian block by block,
which the proof spends through Lemma~\ref{lem:hermitian-constraint-values} ---
with every entry declared \tref{homogeneous}, so that the trusted-mode
right-hand side it denotes is the constructed $c_k\tau_A$
(Definition~\ref{def:homog}). Let $\Feas^{\mathrm{tr}}$ be the
\tref{feasible set} of the trusted-mode declaration and $\Feas^{\mathrm{un}}$
that of the untrusted-mode declaration obtained by replacing the marginal
slot with $\nomarg$ and reading each entry's right-hand side as the scalar
$c_k$, all other slots unchanged. Then
\begin{equation}\label{eq:inclusion}
  \Feas^{\mathrm{tr}}\subseteq\Feas^{\mathrm{un}} ,
\end{equation}
and consequently, for \emph{any} functional $J$ on $\Lin{A\Qsys}$,
$\inf_{\Feas^{\mathrm{un}}}J\le\inf_{\Feas^{\mathrm{tr}}}J$. The inclusion is
\emph{uniform in $\tau_A$}: it holds for every $\tau_A$ satisfying D(ii), and
the untrusted set does not depend on which one the request carried.
\end{proposition}

\begin{proof}
Let $\omega\in\Feas^{\mathrm{tr}}$. The block presentation and $\tr\omega=1$
are shared by the two definitions, as are the per-block positivity conditions;
the untrusted marginal condition is empty, and a set defined by a conjunction
is contained in the set defined by any sub-conjunction, which settles $m=0$
outright and for every $\tau_A$.

For $m\ge1$, fix $k$ and consider first
$\bowtie_k$ equal to $\succeq$, so that
$\tr_{\Qsys}[(\idon{A}\otimes\hat O_k)\omega]\succeq c_k\tau_A$. Both sides are
Hermitian --- the left by Lemma~\ref{lem:hermitian-constraint-values}, using
$\hat O_k$ Hermitian from D(iii) and $\omega$ Hermitian from \eqref{eq:feas},
and the right because $c_k$ is real and $\tau_A$ is Hermitian by D(ii) --- so
the L\"owner comparison is well posed and the traces below are real. Taking
traces of both sides and using that $Y\succeq Y'$ implies $\tr Y\ge\tr Y'$
(the trace of a positive semidefinite operator is non-negative, applied to
$Y-Y'$),
\begin{equation*}
  \tr\!\big[(\idon{A}\otimes\hat O_k)\,\omega\big]
  =\tr\Big(\tr_{\Qsys}\big[(\idon{A}\otimes\hat O_k)\,\omega\big]\Big)
  \;\ge\;\tr(c_k\tau_A)=c_k\,\tr\tau_A=c_k ,
\end{equation*}
using $\tr\tau_A=1$ from D(ii) --- \emph{an identity of the declared numbers,
which is what Definition~\ref{def:exact} delivers and what condition~(S) on the
trace is there to make true} --- and the fact that a partial trace followed by
a trace is the full trace. This is the untrusted entry with $\ge$. The case
$\bowtie_k$ equal to $\preceq$ is identical with the inequality reversed, and
the case $=$ follows by applying both. Hence $\omega\in\Feas^{\mathrm{un}}$,
which is \eqref{eq:inclusion}. The infimum
statement is immediate: an infimum over a superset is not larger than an
infimum over a subset. \emph{Uniformity in $\tau_A$}: the only property of
$\tau_A$ used is $\tr\tau_A=1$, which D(ii) gives for every admissible
marginal, and no $\tau_A$ appears in $\Feas^{\mathrm{un}}$ at all.

\emph{At $d_A=1$ the inclusion is an equality.} There $\tau_A=[1]$ is the only
marginal and $\omega_A=\tau_A$ is the condition $\tr\omega=1$, which
$\Feas^{\mathrm{un}}$ already imposes; and every licensed entry's constructed
right-hand side $c_k[1]$ is the scalar $c_k$ itself, so the two constraint
conditions coincide entry by entry. Hence the two feasible sets are equal and
the two emissions are the same declaration (\xref{transmitter/rem:da1} states
the device-side half of this).
\end{proof}
